% Skriptaufgaben -- MODULNAME (Modul NUMMER), Lektion N
% Quelle: die im Lehrtext-Skript eingebetteten UEBUNGSAUFGABEN (z. B. "Ü 1.3.1")
% am Ende jedes Abschnitts, samt der zugehoerigen offiziellen Loesungen aus dem
% Loesungsteil desselben Skripts. NICHT die Einsendeaufgaben.
%
% Format je Karte:  \lernkarte[id]{Vorderseite}{Rueckseite}
%   Vorderseite = Uebungsnummer (Titelzeile) + Aufgabenstellung (moeglichst
%                 originalgetreu transkribiert -- nur was im Skript steht)
%   Rueckseite  = GANZ KURZ skizzierte Loesung (Kernidee + Ergebnis, 1-3 Zeilen),
%                 abgeleitet aus dem offiziellen Loesungsteil
%
% Diese Datei wird -- wie lernkarten.tex -- NICHT als eigenstaendiges PDF
% kompiliert. generate_docs.py (extract_lernkarten) parst nur die \lernkarte-
% Makros; Vorder- und Rueckseite werden einzeln in Flip-Karten (Web) und
% Anki-Karten gerendert. Der Inhalt beider Argumente ist reines LaTeX -- Mathe
% inline mit $...$ oder abgesetzt mit \[...\]; Aufzaehlungen mit enumitem;
% Tabellen mit array/booktabs.
%
% Verfuegbare Pakete beim Rendern: amsmath, amssymb, enumitem, array, booktabs.
% Kommentarzeilen (mit % am Zeilenanfang) werden ignoriert.
%
% \lernabschnitt{1.1 Reelle Zahlen}  setzt den aktuellen Abschnitt; er wird
% klein ueber den Titel jeder folgenden Karte gesetzt.
%
% \lernkarte[id]{Titel}{Rueckseite}  Die optionale [id] ist ein STABILER,
% inhaltsunabhaengiger Bezeichner; daraus wird die Anki-GUID gebildet. Schema
% fuer Skriptaufgaben: <modul>-<lektion>-sa-<laufnummer>, z. B. [61211-1-sa-01]
% ("sa" = Skriptaufgabe). Titel/Rueckseite duerfen sich spaeter aendern, ohne
% den Anki-Lernfortschritt zu verlieren -- die id NICHT aendern und je Karte
% eindeutig vergeben (neue Karte = neue id). Die id muss global eindeutig sein
% und ist damit disjunkt zu den Notiz-ids <modul>-<lektion>-<nr>.
%
% Im Anki-Deck landen diese Karten unter einem eigenen Unterdeck je Lektion:
%   Mathematikstudium::MODULNAME::Lektion N::Skriptaufgaben

\lernabschnitt{1.1 Titel des Abschnitts}

\lernkarte[MODUL-N-sa-01]{%
  Übungsaufgabe Ü 1.1.1\\[4pt]
  Aufgabenstellung der ersten Übungsaufgabe.
}{%
  Kurze Lösungsskizze -- Kernidee und Ergebnis:
  \[ x = 42 \]
}

\lernkarte[MODUL-N-sa-02]{%
  Übungsaufgabe Ü 1.1.2\\[4pt]
  Aufgabenstellung der zweiten Übungsaufgabe.
}{%
  \begin{itemize}[topsep=2pt,itemsep=1pt]
    \item Schritt / Kernidee
    \item Ergebnis
  \end{itemize}
}
